
\documentclass[11pt]{article}
\usepackage[a4paper,margin=1in]{geometry}
\usepackage{amsmath,amssymb,amsthm}
\usepackage{enumitem}
\usepackage{hyperref}

\newtheorem{theorem}{Theorem}
\newtheorem{proposition}[theorem]{Proposition}
\newtheorem{remark}[theorem]{Remark}
\newtheorem{corollary}[theorem]{Corollary}

\newcommand{\CanSeed}{[2,2,1]/[1,4,4]}

\title{d5\_233: T4 canonical entry from the surfaced canonical prefix and shifted section coordinates}
\author{}
\date{}

\begin{document}
\maketitle

\begin{abstract}
We close the canonical-entry theorem \(T4\) for the surfaced March 21 bundle
semantics that already underlies the closed notes for \(T1,T2,T3\).
At the canonical seed
\[
\CanSeed,
\]
the only live immediate source family after the actual-needed tiny-repair
cleanup is \(R1\).
We show that every such source has the uniform prefix
\[
\partial_1,\ \partial_4,\ \partial_4,
\]
and that after the fixed two-step prefix \(\partial_1\partial_4\), the
standard shifted section coordinates of the explicit one-corner block send the
branch to
\[
E=(m-1,1,u_{\mathrm{source}},2).
\]
This is exactly the manuscript's alternate-\(2\) entry.
\end{abstract}

\section{Setup}

Fix an odd modulus \(m>2\) and the canonical bridge seed
\[
\CanSeed.
\]
In the surfaced endpoint builder this is the candidate with
\[
\texttt{left\_word}=(2,2,1),\qquad \texttt{right\_word}=(1,4,4),
\]
at the representative bridge parameters \(w_0=s_0=0\).

The already closed actual-needed tiny-repair theorem removes the three dead
repair source families \(L3,R2,R3\) at \(\CanSeed\).  So the only remaining
front-end task is to identify the live \(R1\)-continuation.

\section{The exact canonical \(R1\)-family}

\begin{proposition}[canonical \(R1\) source formula]
For color \(0\), the canonical \(R1\)-sources are exactly the vertices
\[
x=(x_0,m-1,0,v,u_{\mathrm{source}})
\]
such that
\[
u_{\mathrm{source}}\neq 0,
\qquad
x_0+v+u_{\mathrm{source}}\equiv 2 \pmod m.
\]
Hence the canonical \(R1\)-family has size \(m(m-1)\).
\end{proposition}

\begin{proof}
In the surfaced builder, color-\(0\) \(R1\)-sources are the layer-\(1\)
vertices in \texttt{layer1\_sets[(0,0)][0]}.  By construction, these are exactly
the vertices with
\[
q=m-1,\qquad w=0,\qquad s=w+u\neq 0,
\qquad x_0+q+w+v+u\equiv 1 \pmod m.
\]
Since \(w=0\) and \(q=m-1\), this is equivalent to
\[
u\neq 0,
\qquad
x_0+v+u\equiv 2 \pmod m.
\]
For each nonzero \(u\) and each \(x_0\), the congruence determines \(v\)
uniquely, so the cardinality is \(m(m-1)\).
\end{proof}

\begin{proposition}[uniform canonical prefix]
Let
\[
x=(x_0,m-1,0,v,u_{\mathrm{source}})
\]
be any canonical \(R1\)-source.
Then the actual continuation under the surfaced canonical seed has the uniform
prefix
\[
x \xrightarrow{\partial_1} x^{(1)} \xrightarrow{\partial_4} x^{(2)}
\xrightarrow{\partial_4} x^{(3)},
\]
where
\begin{align*}
x^{(1)} &= (x_0,0,0,v,u_{\mathrm{source}}),\\
x^{(2)} &= (x_0,0,0,v,u_{\mathrm{source}}+1),\\
x^{(3)} &= (x_0,0,0,v,u_{\mathrm{source}}+2).
\end{align*}
Moreover \(x^{(1)}\) lies in the label class \(R2\), \(x^{(2)}\) lies in the
label class \(R3\), and \(x^{(3)}\) is already background.
\end{proposition}

\begin{proof}
For the canonical seed, the surfaced builder sets the first, second, and third
right directions to
\[
1,\ 4,\ 4.
\]
So every \(R1\)-source first moves by \(\partial_1\), every resulting \(R2\)-site
moves by \(\partial_4\), and every resulting \(R3\)-site again moves by
\(\partial_4\).  Since \(\partial_1\) increments the \(q\)-coordinate and
\(\partial_4\) increments the \(u\)-coordinate, the displayed formulas follow
immediately.
The label statements are exactly how \texttt{\_build\_candidate} constructs the
three right layers.
\end{proof}

\section{Shifted section coordinates and the named entry}

The post-entry one-corner block already uses the surfaced mixed witness reduced
section.  In the surfaced mixed derivation, the reduced first return is written
in the coordinates
\[
(q,s,u,v),\qquad s=w+u,
\]
with exact formula
\[
R(q,s,u,v)=\bigl(q+1,\ s+1+\mathbf 1_{q=m-2},\ u+1,\ v+\mathrm{dv}(q,s)\bigr).
\]
Since \(w=s-u\), this is equivalently the raw section return
\[
R_{\mathrm{sec}}(q,w,u,v)
=
\bigl(q+1,\ w+\mathbf 1_{q=m-2},\ u+1,\ v+\mathrm{dv}(q,w+u)\bigr).
\]

\begin{proposition}[shifted one-corner coordinates]
Define shifted section coordinates by
\[
Q=q-1,\qquad W=w+1,\qquad U=u-1.
\]
Then the surfaced reduced section return becomes
\[
(Q,W,U)\longmapsto
\bigl(Q+1,\ W+\mathbf 1_{Q=m-1},\ U+1\bigr).
\]
This is exactly the section map used by the explicit one-corner odometer block
in the manuscript.
\end{proposition}

\begin{proof}
Substitute \(q=Q+1\), \(w=W-1\), and \(u=U+1\) into the raw section return:
\[
q'=(Q+1)+1,
\qquad
w'=(W-1)+\mathbf 1_{Q=m-1},
\qquad
u'=(U+1)+1.
\]
After shifting back, this becomes
\[
Q'=Q+1,\qquad W'=W+\mathbf 1_{Q=m-1},\qquad U'=U+1.
\]
\end{proof}

\begin{remark}[the suppressed cocycle coordinate]
The front-end entry theorem only needs the \((Q,W,U,\Theta)\)-interface.
The cocycle coordinate \(v\) remains present in the surfaced mixed witness and
reappears later on the graph side, but it is not part of the front-end naming
interface for the one-corner channel.
\end{remark}

\section{T4}

\begin{theorem}[canonical entry theorem \(T4\)]
For every odd modulus \(m>2\), at the canonical bridge seed
\[
\CanSeed,
\]
after the three dead repair source families \(L3,R2,R3\) are removed, the
unique live immediate source family is \(R1\).
Its canonical two-step normalization in the shifted section coordinates is the
alternate-\(2\) entry
\[
E=(m-1,1,u_{\mathrm{source}},2).
\]
\end{theorem}

\begin{proof}
By the canonical support certificate, the only source families present at the
canonical seed are
\[
L3,\ R1,\ R2,\ R3.
\]
By the actual-needed tiny-repair theorem, the families \(L3,R2,R3\) are dead.
So the only live immediate source family is \(R1\).

Now take any canonical \(R1\)-source
\[
x=(x_0,m-1,0,v,u_{\mathrm{source}}).
\]
By the uniform canonical prefix proposition, after the fixed two-step prefix
\(\partial_1\partial_4\) the branch is at
\[
x^{(2)}=(x_0,0,0,v,u_{\mathrm{source}}+1).
\]
On the mixed section, this means
\[
q=0,\qquad w=0,\qquad u=u_{\mathrm{source}}+1.
\]
Applying the shifted coordinates gives
\[
Q=q-1=m-1,\qquad
W=w+1=1,\qquad
U=u-1=u_{\mathrm{source}}.
\]
By convention this is the phase-\(2\) interface of the one-corner block, so the
named entry state is
\[
(Q,W,U,\Theta)=(m-1,1,u_{\mathrm{source}},2)=E.
\]
Since every live canonical source enters the same named interface in this way,
the canonical live continuation is exactly the alternate-\(2\) entry.
\end{proof}

\begin{corollary}[front-end one-corner reduction]
Every unresolved odd-\(m\), \(d=5\) front-end instance reduces to the exact
one-corner channel with canonical seed \(\CanSeed\) and entry state
\[
E=(m-1,1,u_{\mathrm{source}},2).
\]
\end{corollary}

\begin{proof}
The earlier front-end reduction had already compressed every bridge case to the
canonical seed.  Theorem \(T4\) identifies the only remaining live branch there.
\end{proof}

\begin{remark}[saved computational sanity check]
The companion checker
\texttt{d5\_232\_T4\_canonical\_entry\_checker.py} verifies on the exact surfaced
bundle for
\[
m=3,5,7,9,11,13,15,17
\]
that:
\begin{enumerate}[label=\textup{(\arabic*)},leftmargin=2.8em]
\item the canonical \(R1\)-family has size \(m(m-1)\);
\item every source satisfies the explicit source formula;
\item every source has the exact prefix \(\partial_1,\partial_4,\partial_4\);
\item the two-step normalized entry is always
      \(E=(m-1,1,u_{\mathrm{source}},2)\).
\end{enumerate}
This check is only a sanity certificate.  The proof above is symbolic.
\end{remark}

\end{document}
